\chapter{ПРАКТИЧЕСКАЯ РЕАЛИЗАЦИЯ И ТЕСТИРОВАНИЕ}

Во второй главе рассматриваются программная реализация вычислительного ядра,
настольные интерфейсы, сервисный контур, контейнеризация и вопросы верификации.
Основная архитектурная идея проекта состоит в использовании одного ядра для
всех каналов доставки результатов расчёта.

\section{Проектирование структуры программного комплекса}

Программный комплекс реализован как набор взаимосвязанных компонентов,
объединённых общей предметной моделью.
Вычислительное ядро \texttt{if97\_core} выполняет расчёт свойств и не зависит от
конкретного пользовательского интерфейса.
Поверх него построены адаптеры для настольных приложений, web-ориентированного
desktop-варианта и gRPC-сервиса.
Такое разделение позволяет изолировать алгоритмическую часть от кода
представления, обмена данными и развёртывания.

\begin{table}[htbp]
\centering
\caption{Состав программного комплекса}
\begin{tblr}{
colspec={|l|X|},
row{1}={font=\bfseries, halign=c},
hlines,
vlines
}
Компонент & Назначение \\
\texttt{if97\_core} & Вычислительное ядро, реализующее IF97, маршрутизацию расчётов, типы единиц измерения и обработку ошибок \\
\texttt{if97\_app\_api} & Общий слой DTO и контрактов обмена данными между интерфейсами и backend-частями \\
FLTK UI & Нативное настольное приложение для одиночных, табличных и графических расчётов \\
Yew + Tauri & Кроссплатформенный пользовательский интерфейс с web-технологиями и вызовами backend-команд \\
gRPC service & Сервис пакетных расчётов для интеграции во внешнюю инфраструктуру и контейнерную среду \\
\end{tblr}
\end{table}

\subsection{Реализация вычислительного ядра на языке Rust}

Выбор языка Rust обусловлен требованием к надёжности, управляемой
производительности и удобству сопровождения вычислительного кода
\cite{rust_book}.
Ядро организовано послойно.
Слой \texttt{domain} описывает единицы измерения, ошибки и итоговое состояние воды.
Слой \texttt{engine} выполняет валидацию входных данных, определяет регион IF97 и
маршрутизирует запрос к нужному алгоритму.
Слой \texttt{regions} инкапсулирует реализацию региональных уравнений и обратных
зависимостей, а слой \texttt{topology} отвечает за фазовые границы и классификацию
точек расчёта.

Публичный API ядра представлен фасадом \texttt{If97}, предоставляющим методы
\texttt{pt}, \texttt{ph}, \texttt{ps}, \texttt{px}, \texttt{rhot} и
\texttt{metastable\_pt}.
Каждый из этих маршрутов возвращает унифицированную структуру состояния, что
исключает расхождения между интерфейсами и упрощает добавление новых способов
доступа к вычислениям.
Для загрузки коэффициентов IF97 используется build-time генерация таблиц из CSV
в \texttt{build.rs}, что устраняет зависимость от runtime I/O и делает сборку
детерминированной.

\subsection{Реализация настольных приложений и пользовательских сценариев}

Первый пользовательский интерфейс реализован на библиотеке FLTK
\cite{fltk_rs_docs}.
Он ориентирован на быстрый нативный запуск и предоставляет одиночный расчёт,
табличный ввод, управление наборами данных, построение диаграмм и экспорт
результатов.
Особенно важной функцией является отображение купола насыщения и результатов на
термодинамических диаграммах, поскольку именно визуальный анализ часто нужен
при инженерной интерпретации состояния рабочей среды.

Второй интерфейс построен на связке Yew и Tauri \cite{tauri_docs}.
Frontend-часть отвечает за организацию экранов, режимов ввода и отображение
результатов, а Tauri-backend предоставляет команды одиночного расчёта,
табличного расчёта и построения купола насыщения.
Тяжёлые вычисления выполняются в \texttt{spawn\_blocking}, чтобы не блокировать
асинхронный runtime и не ухудшать отзывчивость интерфейса.
За счёт общего слоя DTO сохраняется единый формат обмена между UI и ядром.

\subsection{Контейнеризация, Kubernetes и сервисный контур}

Для интеграции вычислений во внешние системы реализован gRPC-сервис с
двунаправленным потоковым методом расчёта \cite{grpc_docs}.
Формат обмена организован по схеме Struct-of-Arrays, что удобно для пакетной
обработки больших наборов точек и уменьшает издержки сериализации.
Внутри сервиса разделены I/O- и CPU-нагрузки: взаимодействие с сетью выполняет
\texttt{tokio}, а собственно вычисления передаются в пул \texttt{rayon}.
Дополнительно используются ограничения \texttt{in-flight} запросов, что обеспечивает
backpressure и защищает сервис от неконтролируемого роста очередей.

Контейнеризация выполняется с помощью многостадийной сборки Docker
\cite{docker_docs}.
Контейнерный образ содержит только необходимые бинарные артефакты и конфигурацию
сервиса, что уменьшает размер поставки и упрощает развёртывание.
Для эксплуатации в оркестраторе подготовлены манифесты Kubernetes:
Deployment, Service, HPA, PDB и отдельная Job для нагрузочного тестирования
\cite{kubernetes_docs}.
Такой набор позволяет не только запускать сервис, но и описывать его масштабирование,
доступность и поведение под нагрузкой.

\subsection{Автоматизация сборки, тестирование и верификация}

Качество вычислительного ядра подтверждается сочетанием модульных и
интеграционных тестов.
Модульные проверки охватывают региональные зависимости, граничные случаи,
обратные решатели и классификацию фазовой области.
Интеграционные тесты сверяют результаты ядра с контрольными точками из CSV, что
позволяет проверять не отдельные формулы, а весь маршрут расчёта целиком.

Для сервисного слоя дополнительно применяются нагрузочные сценарии, в том числе
через отдельный инструмент \texttt{grpc\_loadtest} и Kubernetes Job.
Автоматизация сборки и публикации организована средствами GitHub Actions
\cite{github_actions_docs}.
Непрерывная интеграция проверяет компиляцию workspace, запуск тестов ядра,
сборку web-части и публикацию документации и релизных артефактов.
В результате получена воспроизводимая цепочка от локальной разработки до
развёртывания сервисного компонента в контейнерной среде.
